\uuid{wjYQ}
\exo7id{6008}
\titre{exo7 6008}
\auteur{quinio}
\organisation{exo7}
\datecreate{2011-05-20}
\isIndication{false}
\isCorrection{true}
\chapitre{Probabilité discrète}
\sousChapitre{Variable aléatoire discrète}
\cpge{mpsi-mp2i.probabilites-univers-fini.c, pcsi.probabilites-variables-aleatoires.c, ptsi.probabilites-variables-aleatoires.c, mpsi-mp2i.probabilites-univers-fini.d, pcsi.probabilites-variables-aleatoires.d, ptsi.probabilites-variables-aleatoires.d, bcpst1.concepts-base-probabilites.a, bcpst1.concepts-base-probabilites.b, bcpst1.variable-aleatoire-univers-fini.d}
\module{Probabilité et statistique}
\niveau{L2}
\difficulte{}

\contenu{
\texte{
L'oral d'un concours comporte au total 100 sujets; les candidats
tirent au sort trois sujets et choisissent alors le sujet traité parmi
ces trois sujets. Un candidat se présente en ayant révisé 60
sujets sur les 100.
}
\begin{enumerate}
    \item \question{Quelle est la probabilité pour que le candidat ait révisé:
}
    \item \question{les trois sujets tirés;}
    \item \question{exactement deux sujets sur les trois sujets;}
    \item \question{aucun des trois sujets.}
\reponse{
Les trois sujets tirés ont été révisés : $P[X=3]=\frac{\binom{60}{3}}{\binom{100}{3}}$.
Deux des trois sujets tirés ont été révisés: $P[X=2]=\frac{\binom{60}{2}.\binom{40}{1}}{\binom{100}{3}}$.
Aucun des trois sujets: $P[X=0]=\frac{\binom{40}{3}}{\binom{100}{3}}$.
}
\end{enumerate}
}
